\chapter{Transient Impact and Propagator Models}\label{ch:mx:transient-impact-and-propagator-models}

The signs of trades are predictable for hours: a buy is followed by buys long after any one trader's order has finished. Yet prices do not trend.
Something must undo the impact of each trade, and how fast it does decides what an execution costs. This chapter measures the memory of order flow in a
simulated market, fits the propagator that maps trades into prices, and asks which shapes of decay a model of impact can have without letting a trader
make money by trading in circles.

\section{The long memory of order flow}

Lillo and Farmer (2004) found that on the London Stock Exchange the autocorrelation of order signs decays roughly as a power law in the lag, with an
exponent of about 0.6: order flow has long memory (One Quant Book 4, chapter~17). Lillo, Mike and Farmer (2005) explained it by splitting: if metaorders
have sizes whose cumulative distribution falls like $v^{-\alpha}$ and are executed in equal pieces, the signs' autocorrelation falls like
$\ell^{-(\alpha-1)}$, a long memory whenever $\alpha<2$.

The chapter's market is chapter~11's \texttt{firm.agentmkt} population, with the noise takers' signs now drawn in runs of Pareto length with tail 1.5, as
if each run were a metaorder. Over two simulated hours, 10\,764 \omterm{def:mx:the-limit-order-book:orders}{market orders} (executions merged), the sign autocorrelation is 0.27 at lag one, 0.106 at
lag ten and 0.023 at lag fifty; a power-law fit over lags 2 to 50 gives an exponent of 0.62 (\cref{fig:mx:transient-impact-and-propagator-models:kernel},
left), against 0.5 for the runs alone: the fundamentalists' orders, whose signs follow the price, dilute the memory. The response
$R(\ell)=\E[\varepsilon_t(p_{t+\ell}-p_t)]$ keeps growing: 0.52 ticks after one order, 1.05 after ten, 2.3 after a hundred.

\section{The propagator model}

\begin{definition}[Propagator model, decay kernel]\label{def:mx:transient-impact-and-propagator-models:propagator}
The \emph{propagator model}\index{propagator model} writes the price before trade $t$ as the sum of the impacts of all past trades,
\[
p_t=\sum_{s<t}G(t-s)\,\varepsilon_s+\text{noise},
\]
where the \emph{decay kernel}\index{decay kernel} (the propagator) $G(\ell)$ is the impact of one trade $\ell$ trades later, independent of the others.
\end{definition}

Bouchaud, Gefen, Potters and Wyart (2004) proposed it and showed that the market sits at a critical point: with correlated signs, a constant $G$ would make
prices trend and a fast-decaying $G$ would make them mean-revert; prices are diffusive only when $G$ decays at the rate that compensates the signs' memory.
The kernel is estimated from the data by writing the price change as $\Delta p_t=\sum_{n\ge0}K(n)\varepsilon_{t-n}$ with $K(n)=G(n+1)-G(n)$ and taking
expectations against past signs: $\E[\varepsilon_{t-k}\Delta p_t]=\sum_nK(n)C(|k-n|)$, a Toeplitz system solved by Levinson recursion.

\omcode{code/firm/propagator/firm_propagator.py}{40}{48}{The kernel from the sign correlations and the sign--price-change cross-correlations, by Levinson recursion.}\label{lst:mx:transient-impact-and-propagator-models:fit}

On the simulated flow the fitted kernel is 0.541 ticks one order after a trade, 0.452 ten orders later and 0.419 a hundred later: it decays, slowly. Does
it make prices diffusive? The variance of price changes over $\ell$ orders divided by $\ell$ is flat for a random walk
(\cref{fig:mx:transient-impact-and-propagator-models:kernel}, right). For the market it falls from 1.98 at lag one to about 1.35 from lag ten on. The
propagator driven by the same signs gives 0.29 at lag one, rising to 1.34 at a hundred: it explains the long-lag variance but not the short-lag noise
(quotes that move without trades), and its decay, cut at a hundred lags, is not quite fast enough. Driven by the same signs shuffled, with no memory, the
same kernel makes prices sub-diffusive, 0.32 falling to 0.18: its decay is what compensates the signs' memory.

\begin{omfigure}
\begin{tikzpicture}
\begin{groupplot}[group style={group size=2 by 1,horizontal sep=1.6cm},width=6.6cm,height=5cm,
  tick label style={font=\footnotesize},label style={font=\small},grid=major,grid style={gray!25},table/col sep=comma,
  xmode=log,xlabel={lag (orders)}]
\nextgroupplot[title={\small signs and kernel},ymode=log,ymin=0.01,ymax=1,
  legend columns=1,legend cell align=left,legend style={font=\footnotesize,at={(0.5,-0.3)},anchor=north,draw=none,fill=none}]
\addplot[omDef,very thick,mark=*] table[x=lag,y=acf]{figdata/microstructure/12-transient-impact-and-propagator-models/kernel.csv};
\addplot[omThm,very thick,mark=square*] table[x=lag,y=G]{figdata/microstructure/12-transient-impact-and-propagator-models/kernel.csv};
\legend{sign autocorrelation $C(\ell)$,kernel $G(\ell)$ (ticks)}
\nextgroupplot[title={\small variance of $p_{t+\ell}-p_t$ over $\ell$},ymin=0,ymax=2.2,
  legend columns=1,legend cell align=left,legend style={font=\footnotesize,at={(0.5,-0.3)},anchor=north,draw=none,fill=none}]
\addplot[black,very thick,mark=*] table[x=lag,y=market]{figdata/microstructure/12-transient-impact-and-propagator-models/diffusion.csv};
\addplot[omThm,very thick,mark=square*] table[x=lag,y=model]{figdata/microstructure/12-transient-impact-and-propagator-models/diffusion.csv};
\addplot[omProp,very thick,mark=triangle*] table[x=lag,y=shuffled]{figdata/microstructure/12-transient-impact-and-propagator-models/diffusion.csv};
\legend{market,propagator on the signs,propagator on shuffled signs}
\end{groupplot}
\end{tikzpicture}

\omcaption{The propagator fitted on two simulated hours. Left: the sign autocorrelation (long memory) and the fitted kernel (slow decay). Right: the
variance of price changes per lag: flat means diffusive. Data: \texttt{mx\_prop.kernel\_study}.}\label{fig:mx:transient-impact-and-propagator-models:kernel}
\end{omfigure}

\section{Transient impact in continuous time}

\begin{definition}[Transient impact model]\label{def:mx:transient-impact-and-propagator-models:transient}
A \emph{transient impact model}\index{transient impact model} prices an execution as the continuous-time propagator: trading at rate $v(s)$ moves the price
at time $t$ by $\int_{s<t}G(t-s)\,v(s)\,ds$, so an execution $x_1,\dots,x_n$ at times $t_1,\dots,t_n$ costs $\tfrac12x^\top\Gamma x$ with
$\Gamma_{ij}=G(|t_i-t_j|)$.
\end{definition}

\begin{definition}[Obizhaeva--Wang model]\label{def:mx:transient-impact-and-propagator-models:ow}
The \emph{Obizhaeva--Wang model}\index{Obizhaeva--Wang model} is the \omterm{def:mx:transient-impact-and-propagator-models:transient}{transient impact model} of a \omterm{def:mx:the-limit-order-book:lob}{limit order book} with a block shape and exponential
resilience: the book refills at rate $\rho$ after being consumed, so $G(t)=e^{-\rho t}$ up to a constant.
\end{definition}

Obizhaeva and Wang (2013) found that the optimal execution in their model mixes discrete and continuous trades: selling $X$ over $[0,T]$, it trades a block
of $X/(\rho T+2)$ at the start, sells at a constant rate $\rho X/(\rho T+2)$, and trades another block of $X/(\rho T+2)$ at the end. The first block uses the
book as it is; the continuous trading matches the book's refilling; the last block uses the refilled book when no later trade will pay for its impact.
Minimising $\tfrac12x^\top\Gamma x$ under $\sum x=X$ gives $x\propto\Gamma^{-1}\mathbf 1$, and on a fine grid with an exponential kernel it reproduces the
two blocks.

\section{No-manipulation conditions}

\begin{definition}[No-dynamic-arbitrage condition]\label{def:mx:transient-impact-and-propagator-models:nda}
An impact model satisfies the \emph{no-dynamic-arbitrage condition}\index{no-dynamic-arbitrage condition} when no round trip, a sequence of buys and sells
with zero net position, has a negative expected cost: $x^\top\Gamma x\ge0$ whenever $\sum x=0$ (the kernel is positive definite on round trips).
\end{definition}

\begin{definition}[Transaction-triggered price manipulation]\label{def:mx:transient-impact-and-propagator-models:ttpm}
\emph{Transaction-triggered price manipulation}\index{transaction-triggered price manipulation} is present when the cheapest way to sell (or buy) a
position includes trades in the opposite direction, although no round trip is profitable.
\end{definition}

Huberman and Stanzl (2004) showed that with time-stationary impact only linear permanent impact rules out quasi-arbitrage; Gatheral (2010) studied which
combinations of impact function and \omterm{def:mx:transient-impact-and-propagator-models:propagator}{decay kernel} allow dynamic arbitrage; Alfonsi, Schied and Slynko (2012) proved that impact must decay as a convex
non-increasing function of time to rule out both standard and transaction-triggered manipulation. The family $G(t)=e^{-(t/\tau)^\kappa}$ shows the
three regimes (\cref{fig:mx:transient-impact-and-propagator-models:family}): convex for $\kappa\le1$; for $1<\kappa\le2$ it is flat, then falls, and is
still positive definite (it is the characteristic function of a stable law); beyond 2 it decays too fast after a flat start and is not.

\omcode{code/firm/propagator/firm_propagator.py}{65}{76}{The most profitable round trip (the smallest eigenvalue on the round-trip subspace) and the optimal liquidation.}\label{lst:mx:transient-impact-and-propagator-models:rt}

On 21 trades over one unit of time with $\tau=0.3$, the best round trip loses 0.136 (in units of the instantaneous impact times the squared size) for
$\kappa=0.5$, 0.042 for $\kappa=1$, 0.008 for 1.5 and 0.002 for 1.8, and earns 0.028 for $\kappa=2.2$, 0.091 for 2.5 and 0.213 for 3. At $\kappa=2$ the
Gaussian kernel's matrix is singular on a fine grid, the boundary itself. For $\kappa=1.5$ and 1.8 no round trip pays but the optimal sale includes two
buys, of $-0.086$ and $-0.33$ of the order: transaction-triggered manipulation. With $\kappa\le1$ it sells only.

\begin{omfigure}
\begin{tikzpicture}
\begin{groupplot}[group style={group size=2 by 1,horizontal sep=1.6cm},width=6.6cm,height=5cm,
  tick label style={font=\footnotesize},label style={font=\small},grid=major,grid style={gray!25},table/col sep=comma,
  xlabel={$\kappa$},xmin=0.3,xmax=3.2]
\nextgroupplot[title={\small best round trip (profit)},ymin=-0.16,ymax=0.24]
\addplot[omDef,very thick,mark=*] table[x=kappa,y=round_trip]{figdata/microstructure/12-transient-impact-and-propagator-models/family.csv};
\draw[gray,dashed] (axis cs:2,-0.16) -- (axis cs:2,0.24);
\nextgroupplot[title={\small smallest trade of the optimal sale},ymin=-0.5,ymax=0.1]
\addplot[omThm,very thick,mark=square*] table[x=kappa,y=min_trade]{figdata/microstructure/12-transient-impact-and-propagator-models/family.csv};
\draw[gray,dashed] (axis cs:1,-0.5) -- (axis cs:1,0.1);
\end{groupplot}
\end{tikzpicture}

\omcaption{Kernels $e^{-(t/0.3)^\kappa}$ on 21 trades. Left: the profit of the best round trip, positive once $\kappa>2$. Right: the smallest trade of the
cheapest sale of one unit, negative (a buy) once $\kappa>1$. Data: \texttt{mx\_prop.kernel\_family}.}\label{fig:mx:transient-impact-and-propagator-models:family}
\end{omfigure}

\section{Estimating the kernel}

The Toeplitz estimator needs only signs and prices, but its answers depend on choices. The maximum lag truncates the kernel, and beyond it the fit treats
impact as permanent. Merging executions into orders, or not, changes what a trade is. Other orders move prices too: quote changes without trades add noise
at short lags, as the simulated market showed. And the model's linearity in signs ignores sizes; extensions let the kernel depend on trade size or on
whether the trade changed the price. The diffusion check is the model's own test: a kernel that leaves prices super- or sub-diffusive is missing some of
the decay.

\section{Tutorial: the kernel that paid you to trade}\label{tut:mx:transient-impact-and-propagator-models:tut}

\begin{tutorial}
\textbf{Goal.} Measure the memory of order flow, fit a propagator, check diffusion, and find the kernels that pay for round trips.
\textbf{End state:} \cref{fig:mx:transient-impact-and-propagator-models:kernel,fig:mx:transient-impact-and-propagator-models:family}.
\begin{tutsteps}
\item \textbf{Flow.} \texttt{mx\_prop.flow(seed, seconds)}: signs and mids of merged \omterm{def:mx:the-limit-order-book:orders}{market orders} from \texttt{firm.agentmkt} with long-memory signs.
\item \textbf{Kernel.} \texttt{firm\_propagator.sign\_acf}, \texttt{response}, \texttt{fit\_kernel}; \texttt{propagate} for the model's prices.
\item \textbf{Diffusion.} Variance ratios for the market, the model and the model on shuffled signs (\texttt{kernel\_study()}).
\item \textbf{Manipulation.} \texttt{impact\_matrix}, \texttt{round\_trip}, \texttt{optimal\_liquidation} over the kernel family (\texttt{kernel\_family()});
\texttt{obizhaeva\_wang(X, rho, T)}; draw with \texttt{fig\_prop.py}.
\end{tutsteps}
\textbf{What to change next.} Fit the kernel separately for orders that moved the price and those that did not; raise the run tail to 1.8 and see the
kernel decay faster.
\end{tutorial}

\section{Build: the propagator}\label{bld:mx:transient-impact-and-propagator-models:propagator}

\begin{build}
\bfield{Purpose} Transient impact for the execution models of chapters 14 and 15, the market-impact simulations of chapter~13, and Book~11's inventory
costs.
\bfield{Interface} \texttt{sign\_acf}, \texttt{response}, \texttt{fit\_kernel(eps, p, maxlag)}, \texttt{propagate(eps, G)}, \texttt{impact\_matrix(times,
kernel)}, \texttt{round\_trip(Gamma)}, \texttt{optimal\_liquidation(Gamma, X)}, \texttt{obizhaeva\_wang(X, rho, T)}.
\bfield{Rules} Time in trades for the discrete model; the Toeplitz system solved by \texttt{scipy.linalg.solve\_toeplitz} (Levinson); round trips on the
subspace $\sum x=0$ with unit norm.
\bfield{Acceptance tests} \texttt{code/firm/propagator/tests/}: a planted power-law kernel recovered from persistent signs; round trips and liquidations
for an exponential, a compressed and a stretched kernel; the discrete optimum close to Obizhaeva--Wang's blocks.
\bfield{Stretch} Size-dependent and price-changing kernels; the continuous-time kernel from timestamps; the non-linear transient impact of Gatheral.
\end{build}

\begin{omsources}
\item J.-P.~Bouchaud, Y.~Gefen, M.~Potters and M.~Wyart, ``Fluctuations and response in financial markets: the subtle nature of `random' price changes'',
\emph{Quantitative Finance} 4(2), 2004.
\item F.~Lillo and J.~D.~Farmer, ``The long memory of the efficient market'', \emph{Studies in Nonlinear Dynamics and Econometrics} 8(3), 2004.
\item G.~Huberman and W.~Stanzl, ``Price manipulation and quasi-arbitrage'', \emph{Econometrica} 72(4), 2004.
\item F.~Lillo, S.~Mike and J.~D.~Farmer, ``Theory for long memory in supply and demand'', \emph{Physical Review E} 71, 2005.
\item J.~Gatheral, ``No-dynamic-arbitrage and market impact'', \emph{Quantitative Finance} 10(7), 2010.
\item A.~Alfonsi, A.~Schied and A.~Slynko, ``Order book resilience, price manipulation, and the positive portfolio problem'', \emph{SIAM Journal on Financial
Mathematics} 3(1), 2012.
\item A.~A.~Obizhaeva and J.~Wang, ``Optimal trading strategy and supply/demand dynamics'', \emph{Journal of Financial Markets} 16(1), 2013.
\end{omsources}

\section{Exercises}

\begin{exercise}[$\star$]\label{exo:mx:transient-impact-and-propagator-models:1}
The sign autocorrelation is 0.27 at lag one. If the next sign were predicted to equal the last, how often would the prediction be right?
\end{exercise}

\begin{exercise}[$\star$]\label{exo:mx:transient-impact-and-propagator-models:2}
Why would a constant kernel make prices trend when signs have long memory?
\end{exercise}

\begin{exercise}[$\star$]\label{exo:mx:transient-impact-and-propagator-models:3}
In the \omterm{def:mx:transient-impact-and-propagator-models:ow}{Obizhaeva--Wang model}, sell 10\,000 shares over two hours with resilience $\rho=4$ per hour. Give the two blocks and the continuous rate.
\end{exercise}

\begin{exercise}[$\star\star$]\label{exo:mx:transient-impact-and-propagator-models:4}
Two trades at times 0 and 1 with $G(t)=e^{-t}$. What does the round trip $(+1,-1)$ cost, and how would you sell one unit with these two trades?
\end{exercise}

\begin{exercise}[$\star\star$]\label{exo:mx:transient-impact-and-propagator-models:5}
Explain why the propagator on shuffled signs makes prices sub-diffusive.
\end{exercise}

\begin{exercise}[$\star\star$]\label{exo:mx:transient-impact-and-propagator-models:6}
Metaorder sizes have a cumulative distribution falling like $v^{-1.5}$ and are executed in equal pieces. What exponent should the sign autocorrelation
have, and why did the simulation find 0.62?
\end{exercise}

\begin{exercise}[$\star\star\star$]\label{exo:mx:transient-impact-and-propagator-models:7}
\emph{Coding.} For $\kappa=2.5$, find the most profitable round trip's trades on the 21-point grid and describe its shape. Why does it earn money?
\end{exercise}

\begin{exercise}[$\star\star\star$]\label{exo:mx:transient-impact-and-propagator-models:8}
\emph{Find the flaw.} ``Our fitted kernel is positive definite, so our execution model is free of manipulation.''
\end{exercise}

\section{Problem: The Kernel That Paid You to Trade}

\begin{problem}[{Weekend problem --- the kernel that paid you to trade}]\label{pb:mx:transient-impact-and-propagator-models:1}
A quant proposes an impact model with a kernel that stays flat for a while and then drops. Check what it implies.

\medskip\noindent\textbf{Part I --- Memory.}
\begin{enumerate}
\item What did Lillo and Farmer find on the London Stock Exchange?
\item State the splitting explanation and its exponent.
\item Give the simulated sign autocorrelation at lags 1, 10 and 50, and its fitted exponent.
\item Give the response after 1, 10 and 100 orders.
\end{enumerate}

\medskip\noindent\textbf{Part II --- The propagator.}
\begin{enumerate}[resume]
\item Write the model and the Toeplitz system for its kernel.
\item Give the fitted kernel at lags 1, 10 and 100.
\item Compare the variance ratios of the market, the model and the model on shuffled signs.
\item Why is the market critical in Bouchaud and co-authors' sense?
\end{enumerate}

\medskip\noindent\textbf{Part III --- Continuous time.}
\begin{enumerate}[resume]
\item Write the cost of an execution in the transient model.
\item Derive the optimal liquidation $x\propto\Gamma^{-1}\mathbf 1$.
\item State the Obizhaeva--Wang solution and explain its two blocks.
\item Solve exercise~3.
\end{enumerate}

\medskip\noindent\textbf{Part IV --- Manipulation and the verdict.}
\begin{enumerate}[resume]
\item Define dynamic arbitrage and transaction-triggered manipulation.
\item State the condition of Alfonsi, Schied and Slynko.
\item Give the best round trip for $\kappa=0.5$, 1, 1.5, 1.8, 2.2, 2.5 and 3.
\item Which kernels make the optimal sale include buys?
\item What happens at $\kappa=2$ exactly?
\item State the \emph{named result}: the round-trip profit under a kernel that decays too fast after a flat start, and the decay parameter at which it vanishes.
\item Which regime is the proposed kernel in?
\item In one sentence: what must a kernel look like to be safe?
\end{enumerate}
\end{problem}

\section{Interview questions}

\begin{interviewq}[$\star$ \iqroles{researcher}]\label{iq:mx:transient-impact-and-propagator-models:1}
Order signs are predictable but prices are not. How is that possible?
\end{interviewq}

\begin{interviewq}[$\star\star$ \iqroles{researcher}]\label{iq:mx:transient-impact-and-propagator-models:2}
How would you estimate the impact kernel from trades and quotes?
\end{interviewq}

\begin{interviewq}[$\star\star$ \iqroles{researcher}]\label{iq:mx:transient-impact-and-propagator-models:3}
What condition must an impact model satisfy to rule out profitable round trips?
\end{interviewq}

\begin{interviewq}[$\star\star$ \iqroles{trader}]\label{iq:mx:transient-impact-and-propagator-models:4}
Why does the optimal execution in the \omterm{def:mx:transient-impact-and-propagator-models:ow}{Obizhaeva--Wang model} start and end with blocks?
\end{interviewq}

\begin{interviewq}[$\star\star$ \iqroles{developer}]\label{iq:mx:transient-impact-and-propagator-models:5}
Compute the cost of a schedule of $n$ trades under a transient kernel. What is the complexity, and how would you make it faster?
\end{interviewq}

\begin{interviewq}[$\star\star\star$ \iqroles{researcher}]\label{iq:mx:transient-impact-and-propagator-models:6}
Your optimiser, fed a fitted kernel, tells a liquidation to buy for a while. Is it a bug?
\end{interviewq}
